% !TeX root = ../navier-stokes-manuscript.tex
\subsection{The trace between adjacent rows}

\begin{lemma}[Adjacent Hilbert-scale trace]\label{lem:AdjacentHilbertScaleTrace}
Let \(s\in\R\) and suppose
\[
  v\in L^2(I_*;H^{s+1}),
  \qquad
  v_t\in L^2(I_*;H^{s-1}).
\]
Then \(v\) has one representative in \(C([0,T_*];H^s)\), and for
\(0\leq a\leq b\leq T_*\),
\begin{equation}\label{eq:AdjacentHilbertScaleEstimate}
  \left|
  \norm{v(b)}_{H^s}^{2}
  -\norm{v(a)}_{H^s}^{2}
  \right|
  \leq
  2\norm{v_t}_{L^2(a,b;H^{s-1})}
  \norm{v}_{L^2(a,b;H^{s+1})}.
\end{equation}
\end{lemma}

\begin{proof}
Let \(A=(1-\Delta)^{1/2}\).  Time mollification and frequency truncation make
the identity
\begin{equation}\label{eq:AdjacentHilbertScaleIdentity}
  \norm{v(b)}_{H^s}^{2}
  -\norm{v(a)}_{H^s}^{2}
  =
  2\operatorname{Re}\int_a^b
  \pair{A^{s-1}v_t(t)}{A^{s+1}v(t)}_{L^2}
  \,dt
\end{equation}
legitimate for smooth approximants.  Cauchy--Schwarz gives
\eqref{eq:AdjacentHilbertScaleEstimate}, and the two assumed \(L^2\) rows
pass the right-hand side to the limit.  The limit identity makes the
\(H^s\) norm continuous.  The same approximation gives weak \(H^s\)
continuity, so the representative is strongly continuous.  Any two
continuous representatives agree everywhere because they agree almost
everywhere.
\end{proof}
